\documentclass[11pt,a4paper]{article}

\usepackage{fontspec}
\usepackage{xeCJK}
\setCJKmainfont{Hiragino Mincho ProN}
\setCJKsansfont{Hiragino Kaku Gothic ProN}
\usepackage{amsmath,amssymb,amsfonts}
\usepackage{graphicx}
\usepackage{hyperref}
\usepackage{bm}
\usepackage{booktabs}
\usepackage{amsthm}
\usepackage[margin=2.5cm]{geometry}
\usepackage{titlesec}

\newtheorem{theorem}{定理}[section]
\newtheorem{proposition}[theorem]{命題}
\newtheorem{corollary}[theorem]{系}
\newtheorem{conjecture}[theorem]{予想}
\newtheorem{lemma}[theorem]{補題}

\newcommand{\Spec}{\mathrm{Spec}}
\newcommand{\Tr}{\mathrm{Tr}}
\newcommand{\Br}{\mathrm{Br}}
\newcommand{\Gal}{\mathrm{Gal}}
\newcommand{\Zp}{\mathbb{Z}[1/p]}
\newcommand{\Qp}{\mathbb{Q}_p}
\newcommand{\Zhat}{\hat{\mathbb{Z}}}
\newcommand{\Fp}{\mathbb{F}_p}
\newcommand{\Gm}{\mathbb{G}_m}

\titleformat{\part}[display]
  {\centering\huge\bfseries}{第 \thepart 部}{20pt}{\Huge}

\begin{document}

\begin{titlepage}
\centering
\vspace*{3cm}
{\Huge\bfseries 算術的真空工学\\[0.5cm]
\Large ボスト-コンヌ相転移からワープドライブへ\\[0.3cm]
--- $\Spec(\mathbb{Z})$ 上の層コホモロジーを経由して ---\\[2cm]}
{\Large Wright Brothers\\[0.5cm]}
{\large 独立研究\\[1cm]}
{\large 2026年4月\\[3cm]}
{\normalsize
全3部構成のモノグラフ：\\[0.5cm]
\textbf{第I部} メカニズムと実験\\
\textbf{第II部} 理論チェーンと未解決問題\\
\textbf{第III部} 層理論的再定式化\\[2cm]
}
\vfill
\end{titlepage}

\tableofcontents
\newpage

% ====================================================================
%  概要
% ====================================================================

\section*{概要}
\addcontentsline{toc}{section}{概要}

本モノグラフは、リーマンゼータ関数の算術構造と、
負の真空エネルギー密度の工学的生成可能性
--- アルクビエレ型ワープ計量の主要要件 ---
を結びつける研究プログラムを提示する。

中心的な観察は、量子真空エネルギーがゼータ関数正則化により計算され、
リーマンゼータ関数がオイラー積
$\zeta(s) = \prod_p (1 - p^{-s})^{-1}$
に分解されるということである。
各素数 $p$ は真空に対する独立な「チャンネル」を寄与する。
素数 $p$ のチャンネルを抑制する操作 ---
\emph{素数チャンネルミュート (prime channel muting)} と呼ぶ ---
は、真空エネルギーに因子 $(1 - p^{-s})$ を乗じるが、
物理的に重要な $s = -3$ ではこの因子は負である。

\textbf{第I部} は、ボスト-コンヌ量子統計力学系を通じたメカニズムを展開し、
符号反転定理を証明し、3つの物理的実装経路を特定し、
超伝導回路を用いた具体的実験を提案する。

\textbf{第II部} は、このメカニズムをBC系からワープドライブに至る
10段階の理論チェーンの中に位置づけ、各段階を確立済み・論証済み・
未解決に分類し、残るスケーリング問題の5つの副問題を定式化する。

\textbf{第III部} は、全体の枠組みを $\Spec(\mathbb{Z})$ 上の
層とコホモロジーの言語で再定式化し、弱エネルギー条件 (WEC) が
離散的な位相的不変量（レギュレータの方向）であることを予想し、
「通常の」真空と「ミュートされた」真空の境界の物理を解析する。

% ====================================================================
%  第I部
% ====================================================================
\newpage
\part{メカニズムと実験}
\label{part:mechanism}

\section{序論}

アルクビエレ・ワープ計量~\cite{Alcubierre1994}は一般相対論の枠内で
超光速移動を許容するが、弱エネルギー条件 (WEC) に違反する物質
--- 負のエネルギー密度 --- を必要とする。
カシミール効果~\cite{Casimir1948}は、量子場への境界条件が
負の真空エネルギーを生み出しうることを実証しているが、
その大きさはアルクビエレ要件に多くの桁で及ばない~\cite{PfenningFord1997}。

本研究の中心的な観察は、量子真空が\emph{算術構造}を持つということである：
そのエネルギーはゼータ関数正則化により計算され、
$\zeta(s)$ は素数上のオイラー積に分解される。
各素数は独立なチャンネルを寄与する。
単一のチャンネルを抑制するだけで真空エネルギーの符号が反転することを示す。

\section{ゼータ正則化真空エネルギー}

\subsection{カシミールエネルギーとゼータ正則化}

モード $\{\omega_n = n\omega_0\}$ を持つ量子場のゼロ点エネルギーは
以下のように正則化される：
\begin{equation}
  E_{\mathrm{vac}}^{\mathrm{reg}} = \frac{\hbar\omega_0}{2}\,\zeta_{\mathcal{H}}(-d),
\end{equation}
ここで $d$ は空間次元、標準的な空洞では $\zeta_{\mathcal{H}}(s) = \zeta(s)$ である。
3次元では：
\begin{equation}
  E_{\mathrm{vac}}^{(3D)} = \frac{\hbar\omega_0}{2}\,\zeta(-3) = \frac{\hbar\omega_0}{240}.
\end{equation}

\subsection{オイラー積と素数チャンネル}

オイラー積
\begin{equation}
  \zeta(s) = \prod_{p\;\text{素数}} \frac{1}{1 - p^{-s}}
  \label{eq:euler}
\end{equation}
は $\zeta(s)$ を独立な素数ごとの寄与に分解する。

\noindent\textbf{定義}（素数チャンネルミュート）.
$p$-ミュートゼータ関数を次で定義する：
\begin{equation}
  \zeta_{\neg p}(s) \equiv \zeta(s)\cdot(1 - p^{-s})
  = \prod_{q \neq p} \frac{1}{1 - q^{-s}}
  = \sum_{\substack{n=1\\p\nmid n}}^{\infty} n^{-s}.
  \label{eq:muted}
\end{equation}
幾何学的には、$\zeta_{\neg p}(s)$ は開部分スキーム
$\Spec(\Zp) = \Spec(\mathbb{Z})\setminus\{(p)\}$ のゼータ関数である。

\section{ボスト-コンヌ系}

ボスト-コンヌ (BC) 系~\cite{BostConnes1995,ConnesMarcolli2006}は
$\ell^2(\mathbb{N}^*)$ 上に作用し、$H|n\rangle = \log(n)|n\rangle$、
分配関数 $Z(\beta) = \zeta(\beta)$ を持つ。
$\beta = 1$ で対称性群
$\Gal(\mathbb{Q}^{\mathrm{ab}}/\mathbb{Q})$
を伴う相転移を示す。

\noindent\textbf{定義}（素数射影）.
$P_{\neg p} = \sum_{p\nmid n} |n\rangle\langle n|$ は
$\mathcal{H}_{\neg p} = \mathrm{span}\{|n\rangle : \gcd(n,p) = 1\}$
への射影である。制限された分配関数は
\begin{equation}
  Z_{\neg p}(\beta) = \Tr_{\mathcal{H}_{\neg p}}(e^{-\beta H})
  = \zeta(\beta)(1 - p^{-\beta}).
\end{equation}
これはオイラー積に由来する厳密な代数的恒等式である。

$P_{\neg p}$ を実現する3つの物理的メカニズム：
\textbf{(A)} スペクトルフィルタリング（$\omega = \log p$ のノッチフィルタ）;
\textbf{(B)} 算術的境界条件（周期 $p$ の周期構造）;
\textbf{(C)} 相転移セクター選択（ガロア作用の $\mathbb{Z}_p^*$ 成分を凍結）.

\section{符号反転定理}

\begin{theorem}[真空エネルギーの符号反転]
\label{thm:sign_flip}
任意の素数 $p \geq 2$ および正の整数 $n$ に対して、
\begin{equation}
  \zeta_{\neg p}(1-2n) = \zeta(1-2n)\cdot(1 - p^{2n-1}).
\end{equation}
$p^{2n-1} \geq 2 > 1$ であるから、因子 $(1-p^{2n-1})$ は狭義に負である。
したがって $\zeta_{\neg p}(1-2n)$ と $\zeta(1-2n)$ は逆符号を持つ。
\end{theorem}

\begin{proof}
$\zeta_{\neg p}(s) = \zeta(s)(1-p^{-s})$ を $s = 1-2n$ で評価。
\end{proof}

\begin{corollary}
$p$-ミュートによる3次元真空エネルギーは
$\zeta_{\neg p}(-3) = (1-p^3)/120 < 0$（全ての $p \geq 2$ で成立）。
\end{corollary}

\begin{table}[h]
\centering
\caption{最初の数個の素数に対する符号反転（$s=-3$）}
\begin{tabular}{@{}ccc@{}}
\toprule
$p$ & $(1-p^3)$ & $\zeta_{\neg p}(-3)$ \\
\midrule
2 & $-7$ & $-7/120$ \\
3 & $-26$ & $-26/120$ \\
5 & $-124$ & $-124/120$ \\
7 & $-342$ & $-342/120$ \\
\bottomrule
\end{tabular}
\end{table}

\section{実験提案}

Nb同一平面導波路共振器（$L = 10\,$mm, $f_0 \approx 6.0\,$GHz）に
10個のSQUIDノッチフィルタ（偶数高調波 $2f_0$ 〜 $20f_0$ に同調）を
組み込んだ $p=2$ 素数フィルタを提案する。
基底温度 $T < 20\,$mK（希釈冷凍機）にて、
熱的占有は無視できる（$\langle n_{\mathrm{th}}\rangle \sim 10^{-9}$）。

\textbf{プロトコル}：全SQUIDを離調した真空エネルギー（構成A）と、
偶数高調波に同調した場合（構成B）を比較。

\textbf{予測}：真空エネルギーの符号反転。
$E_{\mathrm{odd}}/E_{\mathrm{full}} \to \zeta_{\neg 2}(-1)/\zeta(-1) = -1$.

\textbf{成功基準}：
レベル1: $\Delta E \neq 0$;
レベル2: $\zeta$-正則化値と一致;
レベル3: $p=2,3$ で乗法構造を検証;
レベル4: 制御可能な符号反転。

\textbf{必要リソース}：希釈冷凍機を備えた実験室、\$50--100K、12ヶ月。

\section{ワープドライブへの含意}

アルクビエレ計量はバブル壁で $\rho < 0$ を必要とする。
定理~\ref{thm:sign_flip}はまさにこれを生成する。
ただし、回路スケールから巨視的時空曲率への接続には6段階の橋渡しが必要であり、
そのうち最初の3段階のみが本提案で実験的に検証可能である。
$\sim 10^{69}$ 桁のエネルギーギャップ（スケーリング問題）は未解決である。

% ====================================================================
%  第II部
% ====================================================================
\newpage
\part{理論チェーンと未解決問題}
\label{part:chain}

\section{10段階チェーン}

\begin{table}[h]
\centering
\caption{10段階の理論チェーン。E = 確立済み, A = 論証済み, O = 未解決。}
\label{tab:chain}
\small
\begin{tabular}{@{}clll@{}}
\toprule
\# & 段階 & 状態 & 主要文献 \\
\midrule
1 & NCG標準模型（$M^4 \times F$） & E & Connes~\cite{Connes1996} \\
2 & BC系 $=$ $F$ の算術的カーネル & A & Connes-Marcolli~\cite{ConnesMarcolli2006} \\
3 & $\beta = 1$ での相転移 & E & Bost-Connes~\cite{BostConnes1995} \\
4 & 素数ミュート ($P_{\neg p}$) & E & 第I部 \\
5 & 分配関数の恒等式 & E & オイラー積 \\
6 & 真空エネルギー符号反転 & E & 定理~\ref{thm:sign_flip} \\
7 & WEC違反 & E & 標準QFT \\
8 & ワープ計量が許容 & E & Alcubierre~\cite{Alcubierre1994} \\
9 & NCG $\leftrightarrow$ バルク時空 & A & \cite{ConnesMarcolli2006,Maldacena1998} \\
10 & スケーリング ($10^{-22}\to 10^{47}\,$J) & O & --- \\
\bottomrule
\end{tabular}
\end{table}

\section{論証済み段階：精密な予想}

\begin{conjecture}[算術的カーネル]
\label{conj:kernel}
$M^4 \times F$ における内部空間 $F$ は、
$\mathcal{A}_{\mathrm{BC}} = C^*(\mathbb{Q}_+^* \backslash \mathbb{A}_f / \Zhat^*)$
を算術的部分代数として含む。$\mathcal{A}_{\mathrm{BC}}$ の KMS 状態は
$M^4$ 上の量子場の真空セクターに対応する。
\end{conjecture}

\textbf{根拠}：(i) $F$ と $\mathcal{A}_{\mathrm{BC}}$ は共にガロア型の自己同型を持つ
非可換 $C^*$-環である;
(ii) $M^4\times F$ 上のスペクトル作用は $\zeta$ 関数の値を生成し、
BC系の分配関数は $\zeta$ \emph{そのもの}である;
(iii) Connes-Marcolli~\cite{ConnesMarcolli2006} の定理4.1は
BC代数を $\mathbb{Q}$-格子の非可換空間として確立する。
\textbf{不足}：$\mathcal{A}_{\mathrm{BC}} \hookrightarrow \mathcal{A}_F$
の明示的な埋め込み。

段階9（NCG $\leftrightarrow$ バルク曲率）には2つの候補的橋渡しがある：
(A) スペクトル作用の変更, (B) ホログラフィック（AdS/CFT型）。
いずれも予想の段階である。

\section{スケーリング問題：5つの副問題}

\textbf{I. コヒーレント増幅}：$N$ 個の同期系 $\Rightarrow E \propto N^2$？
どのような条件下で？ デコヒーレンス率は？

\textbf{II. トポロジカル保護}：
$K_1(\mathbb{Z}) = \mathbb{Z}/2 \to K_1(\Zp) = \mathbb{Z}/2 \times \mathbb{Z}$
は位相的相転移。物理的に実現可能か？

\textbf{III. 臨界増幅}：BC系の $\beta = 1$ における臨界指数は？
素数ミュートは臨界モードに結合するか？

\textbf{IV. ホログラフィック増幅}：算術的 AdS/CFT？ 中心電荷は？

\textbf{V. ブートストラップ}：プランクスケールのワープバブルの安定性は？
エネルギーコストの半径依存性は？

\section{ロードマップ}

フェーズ1（1--2年）：SQUID実験（段階4--6）。
フェーズ2（2--5年）：層理論 + トポロジカル相（第III部）。
フェーズ3（5--10年）：コヒーレント配列 + 臨界現象。
フェーズ4（10年以上）：巨視的負エネルギー（フェーズ1--3が成功した場合）。

% ====================================================================
%  第III部
% ====================================================================
\newpage
\part{層理論的再定式化}
\label{part:sheaf}

\section{動機：カントール対グロタンディーク}

第I部は $\zeta(s) = \sum n^{-s}$ を点上の和（集合論的、カントール）として扱った。
第III部はこれをスキーム $\Spec(\mathbb{Z})$ のゼータ関数
（層理論的、グロタンディーク）として扱う。

違い：
\emph{集合論的ミュート}は和から項を除去する（量的）；
\emph{層的ミュート}は局所化 $\mathbb{Z} \to \Zp$
により基礎空間自体を変更する（質的）。

\section{$\Spec(\mathbb{Z})$ 上のエネルギー層}

\begin{proposition}[エネルギー層]
$\mathcal{F}(D(n)) = \zeta(s)\cdot\prod_{p|n}(1-p^{-s})$
という割り当てと制限写像
$f \mapsto f\cdot\prod_{p|n,\,p\nmid m}(1-p^{-s})$ は、
$\Spec(\mathbb{Z})$ 上の層（ザリスキ位相）を定義する。
\end{proposition}

大域切断：$\mathcal{F}(\Spec(\mathbb{Z})) = \zeta(s)$。
$D(p)$ 上の切断：$\mathcal{F}(D(p)) = \zeta_{\neg p}(s)$。

\section{局所化完全列}

$\Spec(\mathbb{Z})_{\mathrm{\acute{e}t}}$
上の乗法群層 $\Gm$ に対して：
\begin{equation}
  0 \to \underbrace{\Qp/\mathbb{Z}_p}_{H^2_{(p)}} \to \Br(\mathbb{Z})
  \to \Br(\Zp) \to 0.
\end{equation}

$p$ での局所化はブラウアー群から $\Qp/\mathbb{Z}_p$
（全ての $p$-局所障害）を除去する。

\section{K理論とレギュレータ}

$K_1(\mathbb{Z}) = \mathbb{Z}/2$, $K_1(\Zp) = \mathbb{Z}/2 \times \mathbb{Z}$
（ミュートする素数ごとに自由階数が1増加）。
$K_2(\mathbb{Z}) = \mathbb{Z}/2$, 生成元 $\{-1,-1\}$ (Milnor)。

ボレル・レギュレータ $\mathrm{reg}: K_{2n-1}(\mathbb{Z}) \to \mathbb{R}^{d_n}$
は、余体積が $\zeta(n)$ に関係する格子に写像する~\cite{Borel1977,Lichtenbaum1973}。

\begin{theorem}[レギュレータの符号反転]
任意の素数 $p \geq 2$ および正整数 $n$ に対し、
局所化されたレギュレータの値
$\zeta_{\neg p}(1-2n) = \zeta(1-2n)\cdot(1-p^{2n-1})$
は $\zeta(1-2n)$ と逆符号を持つ。
\end{theorem}

数値的には定理~\ref{thm:sign_flip}と同一だが、K理論的解釈が新しい：

\begin{conjecture}[レギュレータの方向]
\label{conj:reg}
$\zeta(1-2n)$ の符号は位相的不変量である：
ボレル・レギュレータ格子の\emph{方向 (orientation)}。
局所化はこの方向を反転させる。
物理的には、これがWECの成立（正）か違反（負）かを決定する。
\end{conjecture}

正しければ、WECは離散的（オン/オフ）な不変量であり、
連続的なエネルギー閾値ではない。スケーリング問題は解消される。

\section{局所化の5つの物理的解釈}

$\mathbb{Z} \to \Zp$ は5つの等価な物理的読み方を許す：

\begin{enumerate}
\item \textbf{ゲージ化}：$|n\rangle \sim |pn\rangle$ の同一視
  （離散ゲージ対称性）。
\item \textbf{カルツァ-クライン}：$p$-adic余剰次元の非コンパクト化
  （$R_p \to \infty$）。
\item \textbf{相転移}：$p$ の序列パラメータが凍結
  （選択的 $\beta_p \to \infty$）。
\item \textbf{トポス論理}：ハイティング代数が変化；WECの真偽値が反転。
\item \textbf{アデール的}：$\Qp$ 成分を $\mathbb{A}_\mathbb{Q}$ から除去。
\end{enumerate}

実験的に最もアクセスしやすいもの：(1) 量子ビット配列での離散ゲージ理論;
(3) 量子シミュレータでの相転移;
および凝縮系で実現される K 理論的位相相転移
$K_1(\mathbb{Z}) \to K_1(\Zp)$。

\section{境界物理：算術的ドメインウォール}

局所化を有限領域に適用すると、境界
$V(p) = \Spec(\Fp)$ がWEC成立域とWEC違反域を分離する。

局所コホモロジー $H^2_{(p)} \cong \Qp/\mathbb{Z}_p$ が壁上に局在する：
解像度レベル $n$ で $p^n$ 個のモードを持つ無限のタワー
---「算術的エッジ状態」---  である。

\textbf{壁の予測される現象}：
(i) トポロジカルエッジ状態（$\Delta K_1 = \mathbb{Z}$ に由来）；
(ii) $T_{\mathrm{wall}} \sim \hbar\omega_0\log(p)/(2\pi k_B)$ での p-adic 放射；
(iii) 屈折率の不連続 $n_{\mathrm{in}}/n_{\mathrm{out}} = \sqrt{1-1/p}$；
(iv) 穏やかな分岐 (tame ramification) によるトポロジカル安定性
（壁の最小厚 $\sim \ell_P\log p$）。

この算術的ドメインウォールがワープバブルの壁である：
$\Spec(\mathbb{Z})$（通常真空）と $\Spec(\Zp)$（WEC違反真空）を分離する
空間中の閉曲面であり、アルクビエレの形状関数は構造層の遷移に
エンコードされている。

\section{予想の段階にあるもの}

\begin{enumerate}
\item 予想~\ref{conj:reg}（WEC = レギュレータの方向）は未証明。
\item 壁の現象は、数学的バルク-境界対応が物理的内容を持つことを前提としている。
\item スケーリング問題の解消は予想~\ref{conj:reg}に完全に依存する。
  WECが本質的に連続的（エネルギー的）条件であれば、スケーリングは残る。
\end{enumerate}

% ====================================================================
%  結論
% ====================================================================
\newpage
\section*{結論}
\addcontentsline{toc}{part}{結論}

本モノグラフは算術的真空工学の3段階分析を提示した：

\textbf{レベル1（第I部）}：ゼータ正則化真空エネルギーの符号を反転させる
代数的に厳密なメカニズム（BC系による素数チャンネルミュート）と、
具体的かつ実現可能な実験提案。

\textbf{レベル2（第II部）}：このメカニズムをワープドライブに接続する
10段階の理論チェーン。7段階が確立済み、2段階が精密な予想として論証済み、
1段階（スケーリング）が未解決。

\textbf{レベル3（第III部）}：符号反転を $\Spec(\mathbb{Z})$ 上の
レギュレータ方向の反転として再解釈する層理論的再定式化。
エネルギー蓄積の問題ではなく位相的相転移として
スケーリング問題を再構成する可能性を示す。

鍵となる実験的テスト --- SQUID素数フィルタ実験 --- は
既存技術で \$50--100K、12ヶ月で実施可能である。
レベル3の成功（複数素数に対する乗法構造の確認）は、
オイラー積公式が量子真空を支配する物理法則であることを確立し、
ワープドライブへの含意とは独立に、基礎物理学への重要な貢献となる。

\begin{thebibliography}{30}

\bibitem{Alcubierre1994}
M.~Alcubierre, Class. Quantum Grav. \textbf{11}, L73 (1994).

\bibitem{Casimir1948}
H.~B.~G.~Casimir, Proc. K. Ned. Akad. Wet. \textbf{51}, 793 (1948).

\bibitem{PfenningFord1997}
M.~J.~Pfenning and L.~H.~Ford, Class. Quantum Grav. \textbf{14}, 1743 (1997).

\bibitem{BostConnes1995}
J.-B.~Bost and A.~Connes, Selecta Math. (N.S.) \textbf{1}, 411 (1995).

\bibitem{ConnesMarcolli2006}
A.~Connes and M.~Marcolli, \emph{Noncommutative Geometry, Quantum Fields
and Motives}, AMS (2008); Theorem~4.1 and Chapter~3.

\bibitem{Connes1996}
A.~Connes, Commun. Math. Phys. \textbf{182}, 155 (1996).

\bibitem{Wilson2011}
C.~M.~Wilson et al., Nature \textbf{479}, 376 (2011).

\bibitem{Lamoreaux1997}
S.~K.~Lamoreaux, Phys. Rev. Lett. \textbf{78}, 5 (1997).

\bibitem{Maldacena1998}
J.~Maldacena, Adv. Theor. Math. Phys. \textbf{2}, 231 (1998).

\bibitem{Lentz2021}
E.~W.~Lentz, Class. Quantum Grav. \textbf{38}, 075015 (2021).

\bibitem{BobrickMartire2021}
A.~Bobrick and G.~Martire, Class. Quantum Grav. \textbf{38}, 105009 (2021).

\bibitem{Borel1977}
A.~Borel, Ann. Scuola Norm. Sup. Pisa \textbf{4}, 613 (1977).

\bibitem{Lichtenbaum1973}
S.~Lichtenbaum, in \emph{Algebraic $K$-theory II}, Springer LNM \textbf{342}, 489 (1973).

\bibitem{Bordag2009}
M.~Bordag et al., \emph{Advances in the Casimir Effect}, Oxford (2009).

\bibitem{Kitaev2009}
A.~Kitaev, AIP Conf. Proc. \textbf{1134}, 22 (2009).

\bibitem{Quillen1973}
D.~Quillen, in \emph{Algebraic $K$-theory I}, Springer LNM \textbf{341}, 85 (1973).

\end{thebibliography}

\end{document}
